% Part: proof-theory
% Chapter: sequent-calculus
% Section: quantifiers

\documentclass[../../../include/open-logic-section]{subfiles}

\begin{document}

\olfileid[id]{pt}{seq}{qua}

\olsection{Bukti Reguler dan Substitusi}

Aturan-aturan untuk kuantor menimbulkan beberapa kerumitan karena adanya
``syarat variabel eigen'' yang berlaku bagi aturan \LeftR{\lexists}
dan~\RightR{\lforall}. Sebagian besar kerumitan ini dapat ditangani dengan
memastikan bahwa !!{constant}~$a$ dalam inferensi-inferensi tersebut tidak
pernah digunakan kembali. Mari kita perkenalkan definisi berikut.

\begin{defn}
!!^a{proof}~$\pi$ dalam \Log{G1c} atau \Log{G3c} disebut \emph{reguler} jika
\begin{enumerate}
  \item tidak ada variabel eigen~$a$ yang digunakan sebagai variabel eigen
  bagi lebih dari satu inferensi \LeftR{\lexists} atau~\RightR{\lforall}, dan
  \item jika $a$ merupakan variabel eigen bagi suatu inferensi
  \LeftR{\lexists} atau~\RightR{\lforall}, variabel eigen tersebut hanya muncul di atas
  inferensi tersebut dalam~$\pi$.
\end{enumerate}
\end{defn}

\begin{lem}\ollabel{lem:subst-regular} Jika $\pi$ reguler dan berakhir pada
$\Gamma \Sequent \Delta$, $c$ adalah !!a{constant} yang tidak digunakan
sebagai variabel eigen dalam~$\pi$, dan $t$~adalah term yang tidak memuat
variabel eigen apa pun dari~$\pi$, maka $\Subst{\pi}{t}{c}$ yang dihasilkan
dari~$\pi$ dengan mengganti setiap kemunculan~$c$ oleh~$t$ merupakan
!!a{proof} reguler. Sekuen akhir $\Subst{\pi}{t}{c}$ ialah
$\Subst{\Gamma}{t}{c} \Sequent \Subst{\Delta}{t}{c}$, dengan
$\Subst{\Gamma}{t}{c} = \Setabs{!A(t)}{!A(c) \in \Gamma}$.
\end{lem}

\begin{proof}
  Dengan induksi pada $\pheight{\pi}$. Jika $\pheight{\pi} = 0$, bukti itu
  hanya terdiri atas sebuah aksioma. Maka $\Subst{\pi}{t}{c}$ juga
  merupakan aksioma, sebab setiap !!{formula}~$!A(c)$ di dalamnya berubah
  menjadi $!A(t)$.

  Jika $\pheight{\pi} > 0$, kita membedakan kasus menurut inferensi
  terakhir dalam~$\pi$. Kasus untuk aturan struktural (dalam \Log{G1c})
  dan aturan logis bagi penghubung proposisional mudah dan diserahkan
  sebagai latihan. Kita meninjau kasus-kasus ketika $\pi$ berakhir dengan
  inferensi kuantor.
  \begin{enumerate}
    \item Jika inferensi terakhir adalah $\RightR{\lforall}$, $\pi$
    berbentuk
    \[
    \AxiomC{}
    \RightLabel{$\pi'$}
    \Deduce$\Gamma(c) \fCenter \Delta'(c), !A(a,c)$
    \RightLabel{\RightR{\lforall}}
    \UnaryInf$\Gamma(c) \fCenter \Delta'(c), \lforall[x][!A(x,c)]$
    \DisplayProof
    \]
    Menurut hipotesis induksi, $\Subst{\pi'}{t}{c}$ merupakan
    !!a{proof} reguler atas $\Gamma(t) \fCenter \Delta'(t), !A(a,t)$.
    Karena $a$ merupakan variabel eigen dalam~$\pi$, $a$~tidak muncul dalam~$t$.
    Oleh karena itu, inferensi terakhir dalam $\Subst{\pi}{t}{c}$,
    \[
    \AxiomC{}
    \RightLabel{$\Subst{\pi'}{t}{c}$}
    \Deduce$\Gamma(t) \fCenter \Delta'(t), !A(a,t)$
    \RightLabel{\RightR{\lforall}}
    \UnaryInf$\Gamma(t) \fCenter \Delta'(t), \lforall[x][!A(x,t)]$
    \DisplayProof
    \]
    sah: kesimpulannya tidak memuat~$a$.
    \item Jika inferensi terakhir adalah $\LeftR{\lforall}$ dalam
    \Log{G3c}, $\pi$ berbentuk
    \[
    \AxiomC{}
    \RightLabel{$\pi'$}
    \Deduce$!A(s(c),c), \lforall[x][!A(x,c)], \Gamma(c) \fCenter \Delta'(c)$
    \RightLabel{\LeftR{\lforall}}
    \UnaryInf$\lforall[x][!A(x,c)], \Gamma(c) \fCenter \Delta'(c)$
    \DisplayProof
    \]
    Menurut hipotesis induksi, $\Subst{\pi'}{t}{c}$ merupakan
    !!a{proof} reguler atas $!A(s(t),t), \lforall[x][!A(x,t)], \Gamma(t)
    \fCenter \Delta'(t)$. Inferensi terakhir dalam
    $\Subst{\pi}{t}{c}$,
    \[
    \AxiomC{}
    \RightLabel{$\Subst{\pi'}{t}{c}$}
    \Deduce$!A(s(t),t), \lforall[x][!A(x,t)], \Gamma(t) \fCenter \Delta'(t)$
    \RightLabel{\LeftR{\lforall}}
    \UnaryInf$\lforall[x][!A(x,t)], \Gamma(t) \fCenter \Delta'(t)$
    \DisplayProof
    \]
    sah. Kasus untuk \LeftR{\lforall} dalam~\Log{G1c} serupa, tetapi
    sedikit lebih sederhana.
    \item Inferensi terakhir adalah \RightR{\lexists} atau
    \LeftR{\lexists}: latihan.
  \end{enumerate}
  Dalam setiap kasus, kita menambahkan satu sekuen pada akhir sebuah
  !!a{proof} reguler (atau dua !!{proof}s reguler dalam kasus aturan
  berpremis dua). Sekuen baru itu berbeda dari sekuen akhir $\pi$ hanya
  karena $c$ diganti dengan term~$t$. Selain itu, variabel-variabel eigen
  dalam~$\pi$ dan $\Subst{\pi}{t}{c}$ sama. Karena $t$ diandaikan tidak memuat
  variabel eigen dari~$\pi$, sekuen akhir yang baru tidak memuat variabel
  eigen dari $\Subst{\pi}{t}{c}$. Akibatnya, $\Subst{\pi}{t}{c}$ reguler.
\end{proof}

\begin{prop}
Jika $\pi$ merupakan !!a{proof} dalam \Log{G1c} atau \Log{G3c}, maka
terdapat !!a{proof} reguler~$\pi'$ dengan struktur yang sama (khususnya,
!!{height} yang sama) yang berbeda dari $\pi$ hanya karena variabel-variabel
eigennya diganti.
\end{prop}

\begin{proof}
Khusus untuk keperluan bukti ini, sebut suatu inferensi dengan variabel eigen
dalam~$\pi$ \emph{bersih} jika variabel eigennya bukan variabel eigen bagi
inferensi lain dan tidak muncul di mana pun dalam~$\pi$ kecuali dalam
sub-!!{proof} langsung dari inferensi tersebut; jika tidak, sebut inferensi
itu \emph{kotor}. Misalkan $n$ adalah banyaknya inferensi kotor dengan variabel
eigen dalam~$\pi$. Kita menunjukkan, dengan induksi pada~$n$, bahwa
$\pi$~dapat diubah menjadi !!{proof} reguler~$\pi'$ yang diperlukan. Jika
$n=0$, semua inferensi dengan variabel eigen dalam~$\pi$ bersih, sehingga $\pi$
reguler dan tidak ada lagi yang harus dibuktikan. Jika tidak, pilih suatu
inferensi kotor dengan variabel eigen yang posisinya paling tinggi, misalnya
\RightR{\lforall}. Sub-!!{proof}~$\pi_1$ dari~$\pi$ yang berakhir padanya
berbentuk
    \[
    \AxiomC{}
    \RightLabel{$\pi_2$}
    \Deduce$\Gamma \fCenter \Delta, !A(c)$
    \RightLabel{\RightR{\lforall}}
    \UnaryInf$\Gamma \fCenter \Delta, \lforall[x][!A(x)]$
    \DisplayProof
    \]
Perhatikan bahwa karena $c$ merupakan variabel eigen, variabel tersebut tidak
muncul dalam~$\Gamma$, $\Delta$, ataupun $\lforall[x][!A(x)]$. Bukti $\pi_2$
reguler karena tidak memuat inferensi kotor dengan variabel eigen. Ambil
!!a{constant}~$a$ yang tidak muncul dalam~$\pi$. Menurut
\olref{lem:subst-regular}, $\Subst{\pi_2}{a}{c}$ merupakan bukti reguler
atas $\Gamma \fCenter \Delta, !A(a)$, dan kita dapat menerapkan
~$\RightR{\lforall}$ padanya. Jika $\pi_1$ dalam~$\pi$ kita ganti dengan
bukti $\pi_1'$,
    \[
    \AxiomC{}
    \RightLabel{$\Subst{\pi_2}{a}{c}$}
    \Deduce$\Gamma \fCenter \Delta, !A(a)$
    \RightLabel{\RightR{\lforall}}
    \UnaryInf$\Gamma \fCenter \Delta, \lforall[x][!A(x)]$
    \DisplayProof
    \]
kita telah mengganti satu inferensi kotor yang menggunakan variabel eigen dengan inferensi
bersih, dan satu-satunya perbedaan ialah penggantian variabel eigen $c$
dengan~$a$. Bukti yang dihasilkan mempunyai $n-1$ inferensi kotor dengan
variabel eigen, dan hasilnya mengikuti dari hipotesis induksi.
\end{proof}

Kita tidak akan mengacu secara eksplisit pada proposisi yang baru saja
dibuktikan, tetapi kita akan mengandaikan bahwa semua !!{proof}s yang
ditinjau bersifat reguler---jika tidak, kita selalu dapat mengganti
variabel-variabel eigennya agar bukti itu menjadi reguler.

\end{document}
